\environment env_ConfigGlobal
\environment env_Config

\startcomponent Exercises2
\startsection[title={Exercises for Lesson 2}, reference={sec:exercises2}]

Please answer these using the zero-pronoun and が particle (the invisible carriage).
Please answer three times: once as it is actually said, once showing the
zero-pronoun, and once showing what the zero-pronoun represents.

Example:
\startframedtext
Q: I am a cat\crlf
A: ねこだ (or: neko da)\crlf
\hbox to 2em{} =\varnothing{}がねこだ (or: =\varnothing{}ga neko da)\crlf
\hbox to 2em{} =わたしがねこだ (or: =watashi-ga neko da)
\stopframedtext

\startsubsubject[title={Sentences}]

\startitemize[n]
\item I'm Alice
\item It's sunny
\item I drink tea
\item It's White Day
\item I sing a song
\stopitemize

\stopsubsubject

\startsubsubject[title={Vocabulary}]

\startitemize[packed]
\item Sunny: \JPN{はれ} hare
\item Tea: \JPN{おちゃ} ocha {\sl (the \Grapheme{o} is honorific but almost always used)}
\item Drink: \JPN{\ruby{飲}{の}む} nomu
\item Today: \JPN{\ruby{今日}{きょう}} kyou
\item Weather: \JPN{\ruby{天}{てん}\ruby{気}{き}} tenki
\item Song: \JPN{うた} uta
\item White Day: \JPN{ホワイトデー【ほわいとでー】} howaito dē
  {\sl (White Day is a Japanese holiday where men give gifts to women in return
  for the chocolate they received from them on Valentine's Day.)}
\stopitemize

\stopsubsubject

\stopsection
\stopcomponent
